\documentclass[a4paper,11pt,oneside]{memoir}

% Castellano
\usepackage[spanish,es-tabla]{babel}
\selectlanguage{spanish}
\usepackage[utf8]{inputenc}
\usepackage{placeins}

\RequirePackage{booktabs}
\RequirePackage[table]{xcolor}
\RequirePackage{xtab}
\RequirePackage{multirow}

%Citas bibliografía
\usepackage{cite}
\usepackage[numbers,sort]{natbib}

%Tabla
\usepackage{longtable}

% Links
\usepackage[colorlinks]{hyperref}
\hypersetup{
	allcolors = {red}
}

% Ecuaciones
\usepackage{amsmath}

% Euro simbolos
\usepackage{eurosym}

% Rutas de fichero / paquete
\newcommand{\ruta}[1]{{\sffamily #1}}

% Párrafos
\nonzeroparskip


% Imagenes
\usepackage{graphicx}
\newcommand{\imagen}[2]{
	\begin{figure}[!h]
		\centering
		\includegraphics[width=0.9\textwidth]{#1}
		\caption{#2}\label{fig:#1}
	\end{figure}
	\FloatBarrier
}

\newcommand{\imagenflotante}[2]{
	\begin{figure}%[!h]
		\centering
		\includegraphics[width=0.9\textwidth]{#1}
		\caption{#2}\label{fig:#1}
	\end{figure}
}

%COMANDOS GEOINDOOR -- START --

% El comando \imagen_resize nos permite insertar figuras ajustando el tamaño manualmente
\newcommand{\imagenResize}[4]{
	\begin{figure}[h!]
		\centering
		\includegraphics[scale= #1  ]{#2}\label{fig:#4}
		\caption{#3}
	\end{figure}
}

%Comando para añadir espacios de 1cm, como tabulador
\newcommand\tab[1][1cm]{
	\hspace*{#1}
}

%COMANDOS GEOINDOOR -- END --


% El comando \figura nos permite insertar figuras comodamente, y utilizando
% siempre el mismo formato. Los parametros son:
% 1 -> Porcentaje del ancho de página que ocupará la figura (de 0 a 1)
% 2 --> Fichero de la imagen
% 3 --> Texto a pie de imagen
% 4 --> Etiqueta (label) para referencias
% 5 --> Opciones que queramos pasarle al \includegraphics
% 6 --> Opciones de posicionamiento a pasarle a \begin{figure}
\newcommand{\figuraConPosicion}[6]{%
  \setlength{\anchoFloat}{#1\textwidth}%
  \addtolength{\anchoFloat}{-4\fboxsep}%
  \setlength{\anchoFigura}{\anchoFloat}%
  \begin{figure}[#6]
    \begin{center}%
      \Ovalbox{%
        \begin{minipage}{\anchoFloat}%
          \begin{center}%
            \includegraphics[width=\anchoFigura,#5]{#2}%
            \caption{#3}%
            \label{#4}%
          \end{center}%
        \end{minipage}
      }%
    \end{center}%
  \end{figure}%
}

%
% Comando para incluir imágenes en formato apaisado (sin marco).
\newcommand{\figuraApaisadaSinMarco}[5]{%
  \begin{figure}%
    \begin{center}%
    \includegraphics[angle=90,height=#1\textheight,#5]{#2}%
    \caption{#3}%
    \label{#4}%
    \end{center}%
  \end{figure}%
}
% Para las tablas
\newcommand{\otoprule}{\midrule [\heavyrulewidth]}
%

%Table width
\newcolumntype{L}[1]{>{\raggedright\let\newline\\\arraybackslash\hspace{0pt}}m{#1}}

% Nuevo comando para tablas pequeñas (menos de una página).
\newcommand{\tablaSmall}[5]{%
 \begin{table}
  \begin{center}
   \rowcolors {2}{gray!35}{}
   \begin{tabular}{#2}
    \toprule
    #4
    \otoprule
    #5
    \bottomrule
   \end{tabular}
   \caption{#1}
   \label{tabla:#3}
  \end{center}
 \end{table}
}
% GEOINDOOR
% Nuevo comando para tablas pequeñas (menos de una página).
\newcommand{\tablaWithTotal}[6]{%
 \begin{table}[h!]
  \begin{center}
   \rowcolors {2}{gray!35}{}
   \begin{tabular}{#2}
    \toprule
    #4
    \otoprule
    #5
    \bottomrule
    #6
    \bottomrule
   \end{tabular}
   \caption{#1}
   \label{tabla:#3}
  \end{center}
 \end{table}
}

%
%Para el float H de tablaSmallSinColores
\usepackage{float}

%
% Nuevo comando para tablas pequeñas (menos de una página).
\newcommand{\tablaSmallSinColores}[5]{%
 \begin{table}[H]
  \begin{center}
   \begin{tabular}{#2}
    \toprule
    #4
    \otoprule
    #5
    \bottomrule
   \end{tabular}
   \caption{#1}
   \label{tabla:#3}
  \end{center}
 \end{table}
}

\newcommand{\tablaApaisadaSmall}[5]{%
\begin{landscape}
  \begin{table}
   \begin{center}
    \rowcolors {2}{gray!35}{}
    \begin{tabular}{#2}
     \toprule
     #4
     \otoprule
     #5
     \bottomrule
    \end{tabular}
    \caption{#1}
    \label{tabla:#3}
   \end{center}
  \end{table}
\end{landscape}
}

%
% Nuevo comando para tablas grandes con cabecera y filas alternas coloreadas en gris.
\newcommand{\tabla}[6]{%
  \begin{center}
    \tablefirsthead{
      \toprule
      #5
      \otoprule
    }
    \tablehead{
      \multicolumn{#3}{l}{\small continúa desde la página anterior}\\
      \toprule
      #5
      \otoprule
    }
    \tabletail{
      \hline
      \multicolumn{#3}{r}{\small continúa en la página siguiente}\\
    }
    \tablelasttail{
      \hline
    }
    \bottomcaption{#1}
    \rowcolors {2}{gray!35}{}
    \begin{xtabular}{#2}
      #6
      \bottomrule
    \end{xtabular}
    \label{tabla:#4}
  \end{center}
}

%
% Nuevo comando para tablas grandes con cabecera.
\newcommand{\tablaSinColores}[6]{%
  \begin{center}
    \tablefirsthead{
      \toprule
      #5
      \otoprule
    }
    \tablehead{
      \multicolumn{#3}{l}{\small continúa desde la página anterior}\\
      \toprule
      #5
      \otoprule
    }
    \tabletail{
      \hline
      \multicolumn{#3}{r}{\small continúa en la página siguiente}\\
    }
    \tablelasttail{
      \hline
    }
    \bottomcaption{#1}
    \begin{xtabular}{#2}
      #6
      \bottomrule
    \end{xtabular}
    \label{tabla:#4}
  \end{center}
}

%
% Nuevo comando para tablas grandes sin cabecera.
\newcommand{\tablaSinCabecera}[5]{%
  \begin{center}
    \tablefirsthead{
      \toprule
    }
    \tablehead{
      \multicolumn{#3}{l}{\small\sl continúa desde la página anterior}\\
      \hline
    }
    \tabletail{
      \hline
      \multicolumn{#3}{r}{\small\sl continúa en la página siguiente}\\
    }
    \tablelasttail{
      \hline
    }
    \bottomcaption{#1}
  \begin{xtabular}{#2}
    #5
   \bottomrule
  \end{xtabular}
  \label{tabla:#4}
  \end{center}
}



\definecolor{cgoLight}{HTML}{EEEEEE}
\definecolor{cgoExtralight}{HTML}{FFFFFF}

%
% Nuevo comando para tablas grandes sin cabecera.
\newcommand{\tablaSinCabeceraConBandas}[5]{%
  \begin{center}
    \tablefirsthead{
      \toprule
    }
    \tablehead{
      \multicolumn{#3}{l}{\small\sl continúa desde la página anterior}\\
      \hline
    }
    \tabletail{
      \hline
      \multicolumn{#3}{r}{\small\sl continúa en la página siguiente}\\
    }
    \tablelasttail{
      \hline
    }
    \bottomcaption{#1}
    \rowcolors[]{1}{cgoExtralight}{cgoLight}

  \begin{xtabular}{#2}
    #5
   \bottomrule
  \end{xtabular}
  \label{tabla:#4}
  \end{center}
}




\graphicspath{ {./img/} }

% Capítulos
\chapterstyle{bianchi}
\newcommand{\capitulo}[2]{
	\setcounter{chapter}{#1}
	\setcounter{section}{0}
	\chapter*{#2}
	\addcontentsline{toc}{chapter}{#2}
	\markboth{#2}{#2}
}

% Apéndices
\renewcommand{\appendixname}{Apéndice}
\renewcommand*\cftappendixname{\appendixname}

\newcommand{\apendice}[1]{
	%\renewcommand{\thechapter}{A}
	\chapter{#1}
}

\renewcommand*\cftappendixname{\appendixname\ }

% Formato de portada
\makeatletter
\usepackage{xcolor}
\newcommand{\tutor}[1]{\def\@tutor{#1}}
\newcommand{\course}[1]{\def\@course{#1}}
\definecolor{cpardoBox}{HTML}{E6E6FF}
\def\maketitle{
  \null
  \thispagestyle{empty}
  % Cabecera ----------------
\noindent\includegraphics[width=\textwidth]{cabecera}\vspace{1cm}%
  \vfill
  % Título proyecto y escudo informática ----------------
  \colorbox{cpardoBox}{%
    \begin{minipage}{.8\textwidth}
      \vspace{.5cm}\Large
      \begin{center}
      \textbf{TFG del Grado en Ingeniería Informática}\vspace{.6cm}\\
      \textbf{\LARGE\@title{}}
      \end{center}
      \vspace{.2cm}
    \end{minipage}

  }%
  \hfill\begin{minipage}{.20\textwidth}
    \includegraphics[width=\textwidth]{escudoInfor}
  \end{minipage}
  \vfill
  % Datos de alumno, curso y tutores ------------------
  \begin{center}%
  {%
    \noindent\LARGE
    Presentado por  \@author{}\\ 
    en Universidad de Burgos --- \@date{}\\
    Tutor:  \@tutor{}\\
  }%
  \end{center}%
  \null
  \cleardoublepage
  }
\makeatother


% Datos de portada
\title{Geoindoor \\Documentación Técnica}
\author{Juan Pedro Pascual Vitores}
\tutor{Carlos López Nozal}
\date{\today}

\begin{document}

\maketitle



\cleardoublepage



%%%%%%%%%%%%%%%%%%%%%%%%%%%%%%%%%%%%%%%%%%%%%%%%%%%%%%%%%%%%%%%%%%%%%%%%%%%%%%%%%%%%%%%%



\frontmatter


\clearpage

% Indices
\tableofcontents

\clearpage

\listoffigures

\clearpage

\listoftables

\clearpage

\mainmatter

\appendix

\apendice{Plan de Proyecto Software}

\section{Introducción}

La planificación es una parte muy importante del proyecto, ya que es la parte donde se estima el tiempo que se va consumir, el esfuerzo, y si ese proyecto va dar beneficios económicos.

He de decir que a la planificación de los tiempos y del trabajo se hace fundamental si se quiere ser productivo, existen factores no considerados como, que tipo de horas se ha invertido y con que lapsus de tiempo, que afectan a la productividad y a los tiempos. Existen muchos factores que pueden alterar estos parámetros, a continuación citaré algunos.

\begin{itemize}
\item
Horas de desarrollo después de una jornada de trabajo.
\item
Horas de desarrollo dispersas, por ejemplo, se desarrolla durante una hora se para durante horas y se vuelve a retomar, lo que hace que la productividad disminuya .
\item
Curva de aprendizaje, en este proyecto el aprendizaje ha sido fundamental y en muchas ocasiones costoso, si se junta este factor con los 2 anteriores la curva de aprendizaje crece con facilidad.
\item
Horas de desarrollo de poca calidad, es importante tener un lugar donde trabajar donde no se pierda la concentración y el hilo de forma continua.
\end{itemize}

En este anexo se va analizar los factores por los cuales es conveniente o no realizar un proyecto.
\begin{itemize}
\item
Planificación temporal
\item
Viabilidad del proyecto
\item
Viabilidad económica
\item
Viabilidad legal
\end{itemize}
\section{Planificación temporal}

La planificación del tiempo es fundamental, ya que es muchas veces las prisas no nos permiten razonar con serenidad, y tener mucho tiempo lleva a la procrastinación. Entonces lo ideal es llevar un trabajo constante y marcado, es decir crearse horarios de trabajo e intentar cumplirlos lo mejor posible.

Se ha utilizado una metodología ágil Scrum, basada en \emph{Sprints}, y ZenHub, para el seguimiento de issues y tareas. La plataforma GitHub y ZenHub nos permiten administrar el trabajo y las tareas de forma más dinámica y sencilla. Se ha intentado marcar \emph{Sprints} con una fecha final determinada, pero no ha sido posible debido a las dificultades a la hora de estimar el tiempo de trabajo y que no siempre existía la posibilidad de mantener un desarrollo en el tiempo, es decir, no siempre existían horas de calidad para aplicar al desarrollo.
De todas formas si que se han determinado \emph{Sprints} con un inicio y un final, los cuales detallaremos a continuación.
\\
\\
Estos son los pasos que se siguieron en el desarrollo: 

\begin{itemize}
\item
  Aplicar una estrategia de desarrollo incremental a través de
  iteraciones y revisiones.
\item
  La duración media de los \emph{Sprints} debe ser de una semana.
\item
  Al finalizar cada \emph{Sprint} se entrega una parte del producto.
\item
  Se realiza reuniones de revisión al finalizar cada \emph{Sprint} y se vuelve a planificar \emph{Sprint}.
\item
  En la planificación del \emph{Sprint} se genera una pila de tareas.
\item
  Estas tareas se estiman y priorizan en un tablero.
\item
  Para monitorizar el progreso del proyecto se utiliza gráficos
\end{itemize}

Se han seguido todos estos pasos, pero como ya se ha comentado anteriormente la finalización de los \emph{Sprints} no estaba determinada así que su final se tomará cuando se acabaron todos los issues pertenecientes al mismo. Los gráficos en muchas ocasiones no son totalmente representativos del trabajo ni el tiempo utilizado ya que en muchas ocasiones el tiempo invertido no era de calidad, o no se podía invertir tiempo.

\subsection{Sprint 0 (26/3/2017 - 3/4/2017)}\label{splrint0}
El \emph{Sprint} 0 fue el primer \emph{Sprint} que se creo y sirvió de toma de contacto, se afianzaron las ideas y se especificaron los objetivos del proyecto. Se puede decir que ha servido para afianzar ideas sobre el proyecto.
\imagenResize{0.4}{img/anexo/sprint0}{Gráfica Sprint 0}{imgsprint0}


\subsection{Sprint 1 (10/4/2017 - 10/8/2017)}\label{splrint1}
El \emph{Sprint} 1 ha sido uno de los más importantes ya que es donde se ha centrado la mayoría del aprendizaje y una gran cantidad de desarrollo, este \emph{Sprint} ha sido en el que se volcó por primera vez la herramienta Geoindoor en su conjunto, por lo que se crearon commits de gran tamaño.
\imagenResize{0.4}{img/anexo/sprint1}{Gráfica Sprint 1}{imgsprint1}


\subsection{Sprint 2 (10/8/2017 - 30/8/2017)}\label{splrint2}
Si miramos el \emph{Sprint} 2 podemos observar que la fase de aprendizaje está mas avanzada y que el desarrollo se hace menos costoso, pero también esta cargado de menos tareas que el anterior. En este proyecto cabe destacar que el aprendizaje no acaba siempre se encuentra algo nuevo.
\imagenResize{0.4}{img/anexo/sprint2}{Gráfica Sprint 2}{imgsprint2}


\subsection{Sprint 3 (25/9/2017 - 12/10/2017)}\label{splrint3}
El \emph{Sprint} 3 está más dedicado al desarrollo en la parte del viewer, con lo cual se tuvo que volver a investigar y a analizar esa parte de la herramienta para entender cual era la mejor manera para introducir modificaciones. Una ventaja es que entre Architect y Viewer hay un gran parecido, por lo tanto el esfuerzo no fue tanto como la primera vez. También decir que se tuvieron que hacer cambios para que las búsquedas de edificios fueran mas efectivas, lo que llevo un gran esfuerzo.
\imagenResize{0.4}{img/anexo/sprint3}{Gráfica Sprint 3}{imgsprint3}


\subsection{Sprint 4 (17/10/2017 - *)}\label{splrint3}
El \emph{Sprint} 4 es la última iteración antes de la entrega en ella se encuentran los últimos retoques a la herramienta así como, las pruebas realizadas y elementos dedicados a la documentación.
No tiene fecha final ya que abarca, como hemos dicho antes, los últimas pinceladas a la herramienta y la documentación. No encontramos gráfico por no tener fecha final.
 
\section{Estudio de viabilidad}

En cuanto al estudio de la viabilidad, se ha investigado y no hay mucha competencia en cuanto a geolocalización indoor basada en GPS. La mayoría de servicios de geolocalización indoor son de pago, y utilizan bacons para la localización. Lo que hace a Geoindoor uno de los sistemas más atractivos.

\subsection{Viabilidad económica}

Cabe destacar que este sistema en un primer momento no esta pensado para obtener beneficio económico, pero se hará un análisis de la viabilidad económica basándonos en que se vende la herramienta Geoindoor a instituciones e infraestructuras como colegios, institutos, universidades, museos, centros comerciales, palacios de congreso, teatros, etc\ldots.

A continuación haremos un análisis de la viabilidad económica que puede tener Geoindoor.

Para que el sistema Geoindoor sea rentable y fácil de vender se decide asociarse a un experto en marketing y publicidad, que nos ayudará a dar visibilidad al sistema, y nos servirá de contacto con los clientes para obtener beneficio a través de la venta de la herramienta a instituciones y empresas.

En cuanto a los ingresos se estima, a través de un estudio de mercado, que la venta de la herramienta generaría 60137 \euro.

Con esta información podemos analizar la viabilidad económica del proyecto de una forma más detallada:
\\
\\
\textit{*Seguridad social: Seguridad social mensual a cargo de la empresa.}
\\
\textit{*Gasto mensual: Gasto mensual en sueldos.}

\tablaWithTotal{Sueldos y salarios }{llll}{sueldosysalarios}
{ Puesto a desempeñar & Nº Trabajadores & Seguridad social & Gasto mensual\\}{ 
Programador           & 1               & 387\euro         & 1128,75\euro \\
Experto en Marketing  & 1               & 387\euro         & 1128,75\euro \\
}{
\textbf{Total mensuales:}& \textbf{2}   & \textbf{774\euro}& \textbf{22257,50}\euro \\
}

\textit{*Importe: Importe en base imponible.}

\tablaWithTotal{Gatos fijos}{llll}{gastosfijos}
{ Servicios exteriores    & Periodicidad    & Importe    & \%IVA \\}{ 
	Material de oficina                            & Mensual               & 90\euro    & \%0  \\
	Arrendamientos y cánones (alquiler)            & Mensual               & 450\euro   & \%21 \\
	Reparaciones y conservación                    & Mensual               & 100\euro   & \%21 \\
	Servicios de profesionales independientes      & Mensual               & 0\euro     & \%0  \\
	Primas de seguros                              & Mensual               & 118\euro   & \%0  \\
	Publicidad, propaganda y relaciones públicas   & Mensual               & 450\euro   & \%21 \\
	Luz                                            & Mensual               & 175\euro   & \%21 \\
	Teléfono                                       & Mensual               & 124\euro   & \%21 \\
	Internet                                       & Mensual               & 76\euro    & \%21 \\
	Otros gastos                                   & Mensual               & 150\euro   & \%21 \\
}{
	\textbf{Total:}                                &                & \textbf{1733}\euro &  \\
}


\tablaWithTotal{Desglose gastos fijos}{llll}{desglosegastosfijos}
{ Tipo de gasto                    & Mensual         & Anual \\}{ 
	Gastos de personal             & 3031,5\euro     & 36378\euro \\
	Otros gastos de explotación    & 1733\euro       & 20796\euro \\
	Amortización                   & 97\euro         & 1164\euro \\
	Intereses medios de préstamos  & 85\euro         & 1020\euro \\
}{
	\textbf{Total gastos fijos:}& \textbf{4946,5\euro}   & \textbf{59358\euro} \\
}

\tablaSmall{Beneficios - Gastos}{llll}{beneficiosgastos}
{             Año                    & 1º año  \\}{ 
	\textbf{Total beneficios:}       & \textbf{60137\euro} \\
	Gastos personal                  & 36378\euro \\
	Gastos explotación               & 20796\euro \\
	Amortización                     & 1164\euro \\
	Intereses                        & 1020\euro \\   
	\textbf{Total gastos:}           & 59358\euro \\
	Beneficio neto                   & \textbf{779\euro} \\ 
}

Como se puede observar se obtiene un beneficio de 779\euro, con lo cual la conclusión, Geoindoor es un proyecto viable.
 
\subsection{Viabilidad legal}

En cuanto a la viabilidad legal AnyPlace posee una licencia MIT de código abierto lo que nos permite modificar libremente la herramienta. El uso de los mapas de Google Maps y sus facilidades también se encuentra abierto al publico, así que en primera instancia, no nos encontraríamos ningún problema.  Por otra parte hay un tema muy importante a tener en cuenta y es el derecho a la intimidad, y que puede afectar a la ley Orgánica 1/1982, de 5 de mayo, sobre protección civil del derecho al honor, a la intimidad personal y familiar y a la propia imagen.

A lo que nos referimos es que Geoindoor no se hace responsable de la colocación de los planos por la parte de los usuarios, pero se guarda el derecho a eliminar y borrar cualquier cuenta o plano si esta no respeta el derecho a la intimidad de la gente.

Otro tema a tener en cuenta es si se permitiría subir planos de infraestructuras críticas o edificios públicos que se pudieran utilizar con fines dañinos para la sociedad.

Teniendo todos estos factores en cuenta, se concluye con la certeza de que Geoindoor es legal pero, tiene cierto potencial que lo puede convertir en peligroso y se debe controlar.

\apendice{Especificación de Requisitos}

\section{Introducción}
Este anexo recoge la especificación de requisitos, la cual define el comportamiento del sistema desarrollado, es utilizado tanto por el cliente y por el equipo para tener una concepto generalizado de lo que se quiere, para que ambas partes lleguen a una idea común. Para este fin muchas veces son utilizados diagramas los cuales aclaren el concepto y lo que se necesita.

Podemos distinguir entre requisitos funcionales y requisitos no funcionales.
\begin{itemize}
\item
\textbf{Requisito funcional:} es aquel que especifica una función o un servicio que debe cumplir un sistema.
\item
\textbf{Requisito no funcional:} este especifica restricciones sobre el diseño y la implementción del sistema.
\end{itemize}

Las características de una buena especificación de requisitos son definidas por:
\begin{itemize}
\item
\textbf{Completa:} deben estar todos los requerimientos y todas sus relaciones.
\item
\textbf{Consistente:} todos los requerimientos y otros documentos de especificación se debe relacionar de forma coherente.
\item
\textbf{Inequívoca:} se debe ser claro y darse a entender.
\item
\textbf{Correcta:} el sistema debe cumplir con todos los requisitos especificados.
\item
\textbf{Trazable:} los requerimientos deben de estar ordenados y organizados de una manera tal, que sea sencilla su identificación.
\item
\textbf{Priorizable:} los requisitos deben estar organizados de forma jerárquica, según su importancia.
\item
\textbf{Modificable:} los requerimientos deben ser fácilmente modificables.
\item
\textbf{Verificable:} se debe poder probar todo requerimiento.
\end{itemize}


\section{Objetivos generales}
\begin{itemize}
\item
Desarrollar un sistema que permita la localización del individuo en interiores.
\item
Desarrollar un sistema que permita la localización de lugares en interiores.
\item
Desarrollar un sistema que permita el trazado de rutas en interiores.
\item
Desarrollar un sistema que permita mostrar lugares y rutas predefinidas en interiores.
\end{itemize}
\section{Catalogo de requisitos}
A partir de los objetivos anteriormente citados obtenemos los siguientes requisitos.

\subsection{Requisitos Funcionales}

\begin{itemize}
\item
\textbf{RF-1 Gestión de edificio:} la herramienta debe permitir la gestión de un edifico.
\begin{itemize}
	\item
	\textbf{RF-1.1 Emplazamiento de edificio:} la herramienta debe permitir la colocación de un edifico sobre el plano de localización
	\item
	\textbf{RF-1.2 Detallar información del plano:} se debe permitir adjuntar información al edificio (nombre, descripción, etc			\ldots).
	\item
	\textbf{RF-1.3 Emplazamiento de plano:} la herramienta debe permitir la colocación del plano del edifico sobre el plano de 			localización.
	\item
	\textbf{RF-1.4 Redimensión y movimiento del plano:} se debe permitir el movimiento y redimensión del plano añadido.
\end{itemize}
\item
\textbf{RF-2 Gestión de pois:} la herramienta debe permitir la gestión de pois.
\begin{itemize}
	\item
	\textbf{RF-2.1 Emplazamiento de pois:} se debe permitir la colocación de puntos de interés sobre el plano.
	\item
	\textbf{RF-2.2 Detallar información del poi:} se debe permitir adjuntar información al poi (nombre, descripción, tipo, etc			\ldots).
	\item
	\textbf{RF-2.3 Ubicación del poi:} la herramienta debe permitir la colocación del poi sobre una localización del plano de un edifico.
\end{itemize}
\item
\textbf{RF-3 Enlazado de pois:} se debe permitir crear un camino entre pois.
\item
\textbf{RF-4 Gestión de rutas:} se debe permitir la gestión de rutas predefinidas.
\item
\textbf{RF-:}
\item
\textbf{RF-:}
\item
\textbf{RF-:}
\end{itemize}

\section{Especificación de requisitos}



\apendice{Especificación de diseño}

\section{Introducción}
En este apartado, nos detendremos a detallar la estructura de datos, los procedimientos y la arquitectura, que sigue la herramienta Geoindoor.
\section{Diseño de datos}

Este proyecto trabaja con una base de datos no sql por lo tanto la manera en la que están estructurados los datos, es un poco diferente a la que estamos acostumbrados con bases de datos sql.
Los datos en vez de estar estructurados en tablas (BD sql), se estructuran en documentos\footnote{Un documento es un conjunto de entidades de bases de datos que se agrupan en una entidad}. 

Además Geoindoor utiliza el servicio web de bases de datos de Firebase, por lo que la BD esta formada por un único documento. Se usa JSON para el almacenamiento de datos, lo que nos lleva a utilizar un sistema de \textit{clave-valor} para la ordenación.

En conclusión podríamos resumir las características de la BD de Geoindoor en 5 puntos:

\begin{itemize}
	\item Base de datos no sql.
	\item Documentos en vez de entidades relacionales.
	\item Estructura de un único documento.
	\item Almacenamiento de datos a través \textit{clave-valor}
	\item Utilización de JSON para almacenamiento de datos
\end{itemize} 

Entonces la manera de analizar el diseño de datos no se puede realizar desde el punto de vista relacional. De todas formas, a continuación se muestran diagramas que permiten entender mejor el diseño y estructura de datos de Geoindoor.

\imagenResize{0.5}{img/DiagramaRelacionEntidades}{Diagrama de relacion de entidades}{DiagramaRelacionEntidades}

Estos diagramas se leen de izquierda a derecha, para saber que entidad contiene a cual. Las entidades moradas contienen otras entidades, mientras que las azules guardan la relación \textit{clave-valor}, donde se almacenan realmente los datos. Todas las entidades se encuentran en un mismo documento (Geoindoor DB).

La relación virtual indica que hay una relación lógica entre esas entidades, es decir, hay relación de datos entre entidades. En este caso la entidad ``id1 - id..n'', que identifica a los edificios, esta relacionada con ``id1rutas - id..nrutas'', que identifica las rutas de cada edificio. De esta manera podríamos decir que la relación entre ``id1 - id..n'' y ``id1rutas - id..nrutas' es de 1 - *.

\newpage
\textbf{ Geoindoor DB entidades:}
\begin{itemize}
	\item \textbf{Contador}
	\begin{itemize}
		\item Clave 1: id
	\end{itemize}
	\item \textbf{Emails}
	\begin{itemize}
		\item \textbf{e1 - e..n}: identificadores de email.
		\begin{itemize}
			\item Clave 1: contrasena
			\item Clave 2: email 
		\end{itemize}
	\end{itemize}
	\item \textbf{Edificios}
	\begin{itemize}
		\item \textbf{id1 - id..n}: identificadores de edificio.
		\begin{itemize}
			\item Clave 1: idmail 
			\item Clave 2: nombre
		\end{itemize}
		\item \textbf{Rutas}
		\begin{itemize}
			\item \textbf{id1rutas - id..nrutas}: identificadores de ruta de edificio.
			\begin{itemize}
				\item \textbf{r1 - r..rn}: identificadores de ruta. \\
				\hspace{0.5cm} Clave 1: nombre \\
				\hspace{0.5cm} Clave 2: rutas\\
				\hspace{0.5cm} Clave 3: puntos \\
			\end{itemize}
		\end{itemize}
	\end{itemize}
\end{itemize}




\section{Diseño procedimental}

En el diseño procedimental explica como interaccionan las distintas partes del algoritmo.
A continuación se muestran unos diagramas de secuencia que describen el comportamiento de los algoritmos utilizados para el correcto funcionamiento de la herramienta Geoindoor. No se exponen todos los diagramas de secuencia resultantes de el análisis de los algoritmos, por que no son relevantes para el proyecto Geoindoor.

Para los diagramas de la parte de Architect de Geoindoor, el usuario ya esta logeado una vez comienza la acción.
\newpage
\imagenResize{0.5}{img/DiagramasSecuencia/DSCrearEdificio}{Diagrama de secuencia de la creación de un edificio}{DSCrearEdificio}

Para que el siguiente diagrama tenga sentido hay que ubicarle en un contexto. Debe haber un edificio creado y con al menos dos pois.
\newpage
\imagenResize{0.3}{img/DiagramasSecuencia/DSCrearRuta}{Diagrama de secuencia de la creación de una ruta predefinida}{DSCrearRuta}

Para el siguiente diagrama el contexto cambia. Ahora deber haber un edificio con al menos una ruta predefinida.

\imagenResize{0.5}{img/DiagramasSecuencia/DSDibujadoRuta}{Diagrama de secuencia del dibujado de una ruta predefinida al seleccionarla}{DSDibujadoRuta}

\newpage
El siguiente diagrama tiene el mismo contexto que el anterior.

\imagenResize{0.5}{img/DiagramasSecuencia/DSBorrarRuta}{Diagrama de secuencia del borrado de una ruta}{DSBorrarRuta}

\newpage
En este diagrama, nos encontramos en Geoindoor Viewer, en un edificio con al menos una ruta predefinida.

\imagenResize{0.5}{img/DiagramasSecuencia/DSDibujarRutaV}{Diagrama de secuencia del dibujado de una ruta (Viewer)}{DSDibujarRutaV}
 
\section{Diseño arquitectónico}



\apendice{Documentación técnica de programación}

\section{Introducción}

En este apartado se hablará de la estructura final de directorios del proyecto, se detallará información importante para los desarrolladores, como la forma de compilar, instalar y ejecutar el proyecto. Además se explicarán las pruebas realizadas a Geoindoor.

\section{Estructura de directorios}

La estructura de directorios del proyecto en github es la siguiente:

\begin{itemize}
	\item \textbf{./}: directorio raiz donde se encuentra el proyecto Geoindoor.
	\item \textbf{./BaseDeDatos}: directorio donde se encuentra el documento de la exportación de la base de datos.
	\item \textbf{./Despliegue Geoindoor}: directorio donde se encuentra el diagrama del despliegue de la aplicación Geoindoor.
	\item \textbf{./Geoindoor}: directorio donde se encuentra la aplicación Geoindoor pero sin el contexto para el lanzamiento de la herramienta.
	\item  \textbf{./Paquete Geoindoor:} directorio donde se encuentra la aplicación Geoindoor con el contexto necesario para el lanzamiento de la herramienta y el zip que contiene la aplicación Geoindoor.
	\item \textbf{./Planos Museo}: directorio donde se encuentran los planos del museo de la evolución, los cuales serán utilizados para la presentación de Geoindoor.
	\item \textbf{./Server REST API}: directorio donde se encuentra el sistema de ficheros necesario, para crear una REST API en Heroku, la REST API de Geoindoor.
	\footnote{Merece especial mención el fichero \textbf{./Server REST API/index.js} en el cual están las funciones y métodos para la REST API.}
	\item \textbf{./TestEstresGeoindoor}: directorio donde se encuentran los test de estrés de Geoindoor.
	\item \textbf{./TestGeoindoor}: directorio donde se encuentran los test de integración de Geoindoor.
	\item \textbf{./Documentacion}: directorio donde se encuentra la documentación de Geoindoor, escrita en la latex.
\end{itemize}

En cuanto a la estructura de directorios del proyecto se merece especial análisis el \textbf{./Paquete Geoindoor}.

\begin{itemize}
	\item \textbf{Paquete Geoindoor/PaqueteGeoindoor/}: encontramos los lanzadores de la herramienta Geoindoor.
	\item \textbf{Paquete Geoindoor/PaqueteGeoindoor/geoindoor/}: encontramos la página de inicio, con sus respectivas dependencias y las carpetas de las aplicaciones Architect y Viewer.
	\item \textbf{Paquete Geoindoor/PaqueteGeoindoor/geoindoor/architect}: encontramos la página de login y la página de inicio de Architect. Además encontramos las carpetas ``bower\_components'' y ``node\_modules'' donde se encuentran las dependencias de bower y node.
	\item \textbf{Paquete Geoindoor/PaqueteGeoindoor/geoindoor/architect/controllers}: directorio donde se encuentran los controladores de la aplicación.
	\item \textbf{Paquete Geoindoor/PaqueteGeoindoor/geoindoor/architect/scripts}: directorio donde se encuentran los ficheros que contienen las utilidades.\footnote{Se desestima hablar de las restantes carpetas, ya que no son de tanto interés.}
	
	\item \textbf{Paquete Geoindoor/PaqueteGeoindoor/geoindoor/viewer}: encontramos la página de inicio de Viewer. Además encontramos las carpetas ``bower\_components'' y ``node\_modules'' donde se encuentran las dependencias de bower y node.
	\item \textbf{Paquete Geoindoor/PaqueteGeoindoor/geoindoor/viewer/controllers}: directorio donde se encuentran los controladores de la aplicación.
	\item \textbf{Paquete Geoindoor/PaqueteGeoindoor/geoindoor/viewer/scripts}: directorio donde se encuentran los ficheros que contienen las utilidades.\footnote{Se desestima hablar de las restantes carpetas, ya que no son de tanto interés.}
\end{itemize}

\imagenResize{0.3}{img/doctecprog/estructuradirectorios}{Esquema de directorios Geoindoor}{estructuradirectorios}

\section{Manual del programador}

En el proyecto se han utilizado varias herramientas que nos han facilitado el desarrollo, y que habitualmente requieren de instalación.

\begin{itemize}
	\item \textbf{Java} (Tests)
	\item \textbf{Eclipse} (Tests)
	\item \textbf{Git}
	\item \textbf{Node.js}
	\item \textbf{Bower}
	\item \textbf{Grunt}
	\item \textbf{Heroku}
	\item \textbf{Python}
\end{itemize}

\subsection{Java}

Se ha necesitado instalar Java DK 8 (Development Kit) y Java RE 8 (Runtime Environment) para las pruebas de integración a través de selenium. JDK es un software que nos permite crear programas en java \cite{jdkbib}, mientras que JRE son las utilidades que nos permiten correr programas java \cite{jrebib}.

Para la instalación de ambos es, necesario ir a la pagina web oficial (\href{http://www.oracle.com/technetwork/java/javase/downloads/jdk8-downloads-2133151.html}{JDK}, \href{http://www.oracle.com/technetwork/java/javase/downloads/jre8-downloads-2133155.html}{JRE}), ahí debemos buscar la versión que deseemos, y descargarla asegurando que es compatible con nuestro sistema operativo. Una vez descargada, se siguen los pasos del asistente para acabar con la instalación. Posteriormente debemos configurar las variables de entorno \cite{jdkbib}.

\begin{itemize}
	\item \textbf{JAVAPATH}: la ruta completa donde está instalado JDK.
	\item \textbf{CLASSPATH}: la ruta en la que están las bibliotecas o clases de usuario.
	\item \textbf{PATH}: variable en la que se añade la ruta donde esta el JDK.
\end{itemize}


\subsection{Eclipse}
 Para realizar nuestros test de integración con selenium hemos necesitado un IDE (Integrated Development Environment) para realizar la tarea, se ha usado Eclipse Jee Neon \footnote{Eclipse Jee Neon es un IDE que nos permite desarrollar programas java en un entorno determinado, con un contexto determinado.}.
 
 Para la instalación debemos ir a la página oficial (\href{https://www.eclipse.org/downloads/}{Eclipse}) y descargar la versión que queramos, asegurándonos que es compatible con nuestro sistema operativo. Una vez descargado, ejecutamos y seguimos los pasos del asistente hasta finalizar la instalación o en el caso de no haber asistente, se descomprime en el lugar deseado.
 
 \imagenResize{0.25}{img/doctecprog/eclipse}{Eclipse Jee Neon}{eclipse}
 
\subsection{Git} 

Git es un software de control de versiones que ha sido utilizado conjuntamente con GitHub para llevar el control de versiones de Geoindoor. También ha sido utilizado para poner en producción, los cambios en la REST API de Geoindoor (Heroku).

Para instalar Git debemos ir a la página oficial (\href{https://git-scm.com/downloads}{Git}) y descargar la versión que deseemos, fijándonos en que sea compatible con nuestro sistema operativo. Una vez hecha la descarga se sigue el asistente hasta finalizar la instalación. Con la instalación vienen dos herramientas que nos permiten trabajar con Git.

\begin{itemize}
	\item \textbf{Git bash}: consola con la cual se puede trabajar en el control de versiones a través de comandos de Git.
	\item \textbf{Git GUI}: interfaz gráfica que nos permite trabajar con Git de forma más amigable. 
\end{itemize}

\imagenResize{0.5}{img/doctecprog/gitbash}{Terminal de bash de Git}{gitbash}

\imagenResize{0.4}{img/doctecprog/gitgui}{Git GUI}{gitgui}

\imagenResize{0.4}{img/doctecprog/githeroku}{Git en PowerShell}{githeroku}
\textit{Contiene los comandos git utilizados para poner en producción los ficheros en Heroku.}

\subsection{Node.js}

Para poder utilizar bower y grunt temenos que tener instalados el manejador de paquetes ``npm'' y para ello debemos tener instalado Node.js. Para instalar Node.js tenemos que ir a la pagina oficial ()\href{https://nodejs.org/es/download/}{Node.js}) y descargar una versión igual o superior a la 6.11.0 (se recomienda la versión 6.11.0) compatible con nuestro sistema operativo. Una vez se haya descargado, ejecutamos y seguimos los pasos del instalador hasta haber finalizado la instalación.

\imagenResize{0.7}{img/doctecprog/nodenpm}{Comandos ``node'' y ``npm'' una vez instalado Node.js.}{gitgui}

\subsection{Bower}

Bower es un gestor de paquetes para la web, que utilizamos para manejar los paquetes del proyecto, para su instalación, necesita que previamente este instalado node, npm y git. Para la instalación basta con abrir el terminal o consola y ejecutar ``npm install -g bower''.\cite{boweranexbib}

\imagenResize{0.7}{img/doctecprog/bower}{Comando ``bower -v''una vez instalado Bower.}{bower}

\subsection{Grunt}

Grunt es un corredor de tareas (Task runner), que nos ha permitido compilar el proyecto de forma sencilla y rápida, ayudándonos a tener siempre centralizados los scripts y sus dependencias. Nos permite hacer cambios en caliente ya que recompila el proyecto para poder probarlo.

Para la instalación de grunt es necesario tener previamente instalado npm (Node.js). Para instalarlo debemos ejecutar el comando ``npm install -g grunt-cli'' \cite{gruntbibanex}.

\imagenResize{0.7}{img/doctecprog/grunt}{Comando ``grunt''una vez instalado Grunt.}{grunt}

\subsection{Heroku}

Heroku es el servicio web que hemos utilizado para crear la REST API de Geoindoor en Node.js, para poder trabajar con ello debemos descargarnos ``Heroku CLI'' para poder comunicarnos con ser servicio.

Para ello debemos descargarnos e instalar la versión que corresponda a nuestro sistema operativo, desde la página oficial (\href{https://devcenter.heroku.com/articles/getting-started-with-nodejs#set-up}{Heroku CLI}), una vez hecha la descarga, se ejecuta el fichero y se siguen los pasos del instalador hasta finalizar la instalación.

\subsection{Heroku}

Python es un lenguaje de programación interpretado que ha sido utilizado para simular un servidor http, para el desarrollo de Geoindoor no se hace necesario, pero si ofrece comodidad.

Para su instalación vamos a su página oficial (\href{https://www.python.org/downloads/}{Python}) y nos descargamos la versión (mayor o igual a 3.6.1) que corresponda a nuestro sistema operativo. A continuación ejecutamos el fichero descargado y seguimos los pasos del auto instalador hasta finalizar la instalación. En la mayoría de distros de Linux, Python viene ya instalado.

\imagenResize{0.7}{img/doctecprog/pythonan}{Comando ``python -V''una vez instalado Python.}{pythonan}


\section{Compilación, instalación y ejecución del proyecto}
En este apartado se va explicar como compilar, instalar y ejecutar el proyecto Geoindoor desde el principio. Es decir se explicara como se creo la REST API en Heroku de Geoindoor y como se creo la base de datos en Firebase.

\imagenResize{0.35}{img/doctecprog/DespliegueGeoindoor}{Diagrama de despliegue Geoindoor.}{DespliegueGeoindoor}

\subsection{Heroku Server REST API y Firebase Base de datos}

Lo primero que debemos hacer es clonar o descargarnos el repositorio de Geoindoor desde GitHub.


Una vez descargado el repositorio Geoindoor debemos ir a la página oficial de \href{https://id.heroku.com/login}{Heroku} y hacernos una cuenta. Una vez hecha, debemos crearnos un una app en Node.js en Heroku, que se hace mediante los siguientes comandos

\begin{itemize}
	\item \textbf{\textit{heroku login}}: login con la cuenta de Heroku.
	\item \textbf{\textit{git clone https://github.com/heroku/node-js-getting-started.git}}: clonación de la plantilla.
	\item \textbf{\textit{cd node-js-getting-started}}
	\item \textbf{\textit{heroku create}}: creación de la app.
	\item \textbf{\textit{git push heroku master}}: poner la app en el repositorio.
	\item \textbf{\textit{heroku ps:scale web=1}}: asegurarnos de que esta funcionando al menos una instancia de la app.
\end{itemize}

Una vez hecho esto basta sustituir los ficheros del la app recién creada con los ficheros \textbf{./Server REST API/} del repositorio de Geoindoor.

\imagenResize{0.4}{img/doctecprog/herokudash}{Heroku dashboard.}{herokudash}

Ahora para crear el servicio de base de datos en \href{http://firebase.google.com/}{Firebase}, debemos ir a la página oficial y crearnos una cuenta (a través de la cuenta de gmail), desde la página de inicio de la consola creamos un nuevo proyecto, con el nombre que deseemos. Ahora en el apartado de ``Database'' en el panel de control, importamos la base de datos que encontramos en \textbf{./BaseDeDatos/}. 



En este momento debemos enlazar nuestra REST API con la base de datos de Firebase, para ello pegamos en el index.js de nuestra REST API la siguiente porción de código con la información correspondiente, o cambiamos la que está por la correcta.


\begin{verbatim}
firebase.initializeApp({
	serviceAccount: "nombre_del_fichero_de_nuestra_clave_Sin_extension",
	databaseURL: "https://geoindoordb.firebaseio.com"
});
\end{verbatim}
\imagenResize{0.3}{img/doctecprog/paginifire}{Página de inicio de la consola Firebase.}{paginifire}

\imagenResize{0.29}{img/doctecprog/pcontrolfire}{Panel de control Firebase.}{pcontrolfire}

Para generar la clave (serviceAccount) debemos ir al panel de control de nuestra aplicación web en Firebase, y desde Configuración/Cuentas de servicio debemos crear la clave. Una vez creada la ubicamos en la raiz de nuestra REST API y borramos la anterior.

De esta forma quedaría enlazada la REST API con la base de datos de Firebase.

\cite{herokutut}

\subsection{Geoindoor cliente}

Ahora vamos a analizar como compilar, instalar y ejecutar la parte del cliente de Geoindoor.

En primer lugar debemos poner \textbf{./Paquete Geoindoor} en el directorio en el que se quiera trabajar, para lanzar la herramienta se debe ejecutar el lanzador (inicioWindows.bat o inicioLinux.sh), Geoindoor se lanzará en el puerto 9000 (localhost). En Windows, a través de Google Chrome y en Linux a través de Firefox.

Si no se dispone de python instalado, puede usar mismamente \textbf{./Geoindoor} y ubicarlo en el lugar en el que se quiera trabajar, que proporcione un servicio http, ya sea XAMPP, WAMP u otra herramienta que simule un servidor.

Tanto em Viewer como en Architect la forma de comenzar el desarrollo es la misma.

\begin{itemize}
	\item Abirmos la consolola y nos colocamos en la raíz de las aplicaciones, las carpetas \textbf{*/architect/} y \textbf{*/viewer/}
	\item Ejecutamos \textbf{bower install}, si no existiera el fichero ``bower.json''.
	\item Ejecutamos \textbf{grunt}.
\end{itemize}

Si nuestro servicio http funciona bien, nos dirigimos con el navegador a las raíz de la carpeta \textbf{*/architect/} y \textbf{*/viewer/} para comprobar que Geoindoor esta en funcionamiento.
Si aún así no funciona, podemos dirigirnos manualmete a \textbf{*/architect/index.html} y \textbf{*/viewer/index.html} 

\imagenResize{0.9}{img/doctecprog/localhost}{Url de Geoindoor al lanzarla.}{localhost}

\section{Pruebas del sistema}

La pruebas realizadas al sistema Geoindoor, se pueden dividir en dos.

\begin{itemize}
	\item \textbf{Pruebas de integración.}
	\item \textbf{Pruebas de estrés.}
\end{itemize} 

\subsection{Pruebas de integración}

Para las pruebas de integración se ha utilizado Selenium, que es un automatizador de los navegadores, lo que nos permite automatizar acciones sobre la herramienta Geoindoor y ver los resultados.

Las pruebas de integración que se han hecho a Geoindoor son las siguientes.

\tablaSinColores{Requisitos por test de interfaz.}
{L{4cm} L{8cm}}
{2}
{Requisitos por test de interfaz.}
{\textbf{Test} & \textbf{Requisito} \\}
{Gestión de edificios 		&
	\begin{itemize}
		\item RF-1.1 Emplazar edificio.
		\item RF-1.2 Detallar información del edificio.
		\item RF-1.3 Emplazar plano.
		\item RF-1.4 Redimensionar y mover el plano
	\end{itemize}  				
	\\
	Gestión de pois &
	\begin{itemize}
		\item RF-2.1 Emplazamiento de pois.
		\item RF-2.2 Detallar información del poi.
		\item RF-2.3 Ubicar poi.
		\item RF-2.4 Enlazar pois.
	\end{itemize}
	\\  				
	Gestión de rutas &
	\begin{itemize}
		\item RF-3.1 Crear ruta.
		\item RF-3.2 Detallar información de la ruta.
		\item RF-3.3 Añadir pois a ruta.
		\item RF-3.4 Trazar ruta.
		\item RF-3.5 Seleccionar ruta.
		\item RF-3.6 Limpiar ruta.
		\item RF-3.7 Borrar ruta.
	\end{itemize}  				
	\\}

Para la ejecución del test, debemos importar en eclipse el proyecto ``Test'' que se encuentra en \textbf{./TestGeoindoor/Test.zip}. Una vez importado, debemos asegurarnos que tenemos importadas las librerias de Selenium, si no es así debemos hacerlo. Una vez importadas las librerías tenemos que ubicar el driver (Google Chrome, Mozilla Firefox, Safari, etc\ldots), en el lugar que deseemos, sin olvidar que después hay que indicar la ruta del driver en el código de nuestro test. Una vez realizados todos estos pasos, solo nos quedaría lanzar Geoindoor a través de sus lanzadores (inicioLinux.sh o inicioWindows.sh), y ejecutar el test desde eclipse.

\imagenResize{0.35}{img/doctecprog/driverfire}{Test lanzado en Firefox a través de Selenium.}{driverfire}


Los test quedan recogidos en \textbf{./TestGeoindoor/Test.mp4} y el código en java, para el test con Selenium, en \textbf{./TestGeoindoor/Test.zip}. Se proporciona también las librerías de Selenium (\textbf{./TestGeoindoor/selenium-java-3.7.1.zip}) y los drivers de Firefox (\textbf{./TestGeoindoor/geckodriver-v0.19.1-win64.zip}) y Google Chrome (\textbf{./TestGeoindoor/chromedriver\_win32.zip}). Además de la cuenta en Geoindoor para tests (\textbf{./TestGeoindoor/TestCuenta.txt})


\subsection{Pruebas de estrés}

Las pruebas de estrés han sido realizadas a través de la herramienta Webserver Stress Tool, una aplicación de prueba HTTP-client/server, que detecta problemas de rendimiento en un sitio web o servidor web \cite{webserverstresstool}. 

La prueba de estrés consiste en simular 10 usuarios simultáneos, 5 clicks por usuario y 12 segundos entre clicks.


\imagenResize{0.7}{img/doctecprog/tabusers}{Resultados por usuario.}{tabusers}

\imagenResize{0.6}{img/doctecprog/serveruserband}{Ancho de banda del servidor y del usuario.}{serveruserband}

Para ejecutar el test debemos instalar Webserver Stress Tool, una vez instalado, abrimos Webserver Stress Tool y en el apartado de ``Test type'' elegimos la casilla ``CLICKS'', a continuación en el campo ``Run until'' ponemos el valor \textit{5}. En el apartado de ``User simulation'' en ``Number Of Users'', ponemos \textit{10} y en ``Click delay'' \textit{12}. \footnote{En la versión de pago no hay limitación de valores en estos campos, pero en la versión gratuita sí.}

Ahora en el apartado ``URLs'' ponemos la url de Geoindoor, en nuestro caso \textit{http://localhost:9000}. Navegamos hasta el apartado ``Browser Settings'' y en ``Browser Simulation / Use agent'', marcamos la casilla, y ponemos \textit{Mozilla/5.0 (compatible; Webserver Stress Tool 8; Windows)}, en el caso que queramos hacer el test con Firefox en Windows, si no, se debe poner el correspondiente al navegador y sistema operativo. A continuación marcamos la casilla ``Use Timeouts'' y le ponemos un valor de \textit{120}. De esta manera se replica el test que se ha hecho a Geoindoor.

\imagenResize{0.35}{img/doctecprog/webstrestool}{Apartado ``Test Type'' en Webserver Stress Tool.}{webstrestool}


En \textbf{./TestEstresGeoindoor/GeoindoorWebStress.\{pdf doc\}} se encuentra un informe más detallado de los resultados del test, así como la herramienta Webserver Stress Tool (\textbf{./TestEstresGeoindoor/webstress.zip}). Además en \textbf{./TestEstresGeoindoor/TestEstres.\{pdf tex\}} se encuentran las restricciones en cuanto a rendimiento de Firebase y Heroku.








\apendice{Documentación de usuario}

\section{Introducción}

En este apartado hablaremos de cuales son los requisitos del usuario para poder utilizar Geoindoor, los requerimientos para su instalación, y se explicará cómo un usuario puede utilizar la herramienta Geoindoor. 

\section{Requisitos de usuarios}

Los requisitos de un usuario para poder utilizar la aplicación Geoindoor, son los siguientes.

\begin{itemize}
	\item \textbf{Dispositivo con acceso a internet.}
	\item \textbf{Navegador web con JavaScript hablilitado}: se recomienda Firefox o Google Chrome
	\item \textbf{Python 3.6.1 o una versión mayor instalada}: en los distros Linux, la versión no es relevante siempre y cuando tengan el módulo SimpleHTTPServer.
	\item \textbf{El dispositivo debe tener libre el puerto 9000.}
\end{itemize}

\section{Instalación}


En cuanto a la instalación, Geoindoor no necesita una instalación al uso. Lo primero que tenemos que hacer es descargarnos el PaqueteGeoindoor.zip (\textbf{./Paquete Geoindoor/PaqueteGeoindoor.zip}) ya sea directamente desde el repositorio \footnote{Repositorio: jppasvit/geo-indoor} de \href{https://github.com/jppasvit/geo-indoor/tree/master/Paquete\%20Geoindoor}{GitHub}, o desde \href{https://github.com/jppasvit/geo-indoor/releases}{Release} del repositorio. Una vez hecha la descarga, descomprimimos en la ruta en la que queramos que se haga la instalación. Finalizada la descompresión en el directorio deseado, la instalación se da por terminada. Solo queda lanzar la aplicación a través de los lanzadores (inicioWindows.bat y inicioLinux.sh).

\imagenResize{0.65}{img/manualusuario/descomgeo}{Descompresión de PaqueteGeoindoor.zip.}{descomgeo}


\section{Manual del usuario}

En este apartado se darán una serie de indicaciones que facilitará a los usuarios finales el uso de la herramienta Geoindoor.


\subsection{Architect}

Ésta herramienta que forma parte de Geoindoor, nos permitirá crearnos nuestros propios edificios, para que posteriormente en la herramienta Viewer, todo el mundo pueda verlos.

\imagenResize{0.34}{img/manualusuario/pagini}{Página de inicio Geoindoor.}{descomgeo}

\subsubsection{Login}

En primer lugar para acceder a la herramienta Architect, en la pagina de inicio de Geoindoor pulsamos el botón ``Architect''. Esto nos redirigirá a la pagina de Architect, si ya nos habíamos registrado con anterioridad, no hará falta un nuevo login, de manera contraría debemos pulsar en botón login y registrarnos con nuestra cuenta de Gmail. \footnote{Es posible que la aplicación también pida registro para AnyPlace. El usuario debe registrarse también en AnyPlace para poder trabajar.}

\imagenResize{0.52}{img/manualusuario/architect/login}{Página de registro Architect.}{login}

\imagenResize{0.52}{img/manualusuario/architect/regarch}{Registro en Architect con Gmail.}{regarch}

\subsubsection{Crear edificio}

Una vez hecho el registro ya nos encontraríamos en la herramienta Architect sin ninguna restricción.

Para crear un edificio debemos ir a la pestaña ``Buildings'' de nuestro panel de control, pinchar en el icono de ``drag to add new building'' y arrastrar a una posición del plano que nosotros deseemos.

\imagenResize{0.65}{img/manualusuario/architect/pctrlbuild}{Panel de control en la pestaña ``Buildings''.}{pctrlbuild}

Una vez ubicado el edificio aparecerá un panel en el cual podemos decidir el nombre, el código y la descripción del del edificio y si queremos que todo el mundo vea nuestro edificio en la herramienta Viewer.

Una vez rellenada la información correspondiente, se pulsa en el botón ``Add'' para añadir el edificio de forma definitiva.

\imagenResize{0.7}{img/manualusuario/architect/pemergenteedif}{Panel emergente de un edificio.}{pemergenteedif}

Para eliminar un edificio basta con ir a la pestaña ``Delete'' del apartado ``Building Toolbox'' de la pestaña ``Buildings'' de nuestro panel de control, y confirmar el borrado, una vez se ha seleccionado el edificio.\footnote{Se selecciona un edificio pinchando en el icono del rascacielos que aparece en el mapa y que corresponde con su posición, o a través de el desplegable de la pestaña ``Buildings''.}

\subsubsection{Añadir una planta (plano)}

Una vez creado un edificio procedemos a agregarle una planta, un plano de la planta del edificio. Para ello debemos seleccionar el edificio al que queremos añadirle una planta, pinchando sobre el icono del edificio o eligiéndolo a través del desplegable en la pestaña ``Buildings'' del panel de control.

\imagenResize{0.7}{img/manualusuario/architect/iconoedif}{Icono de edificio en el mapa.}{iconoedif}

A continuación, una vez seleccionado el edificio, en nuestro panel de control, pinchamos en la pestaña de ``Floors''. Allí encontraremos un desplegable donde elegir las plantas que ya tenemos creadas, y en el apartado ``Floor Toolbox'', en la pestaña ``Add'', elegimos el numero de planta y damos al botón ``Seleccionar archivo'' para subir el plano de la planta.

Al subir el plano de la planta, la imagen del plano tendrá las típicas herramientas de edición de imagen para girarla y dimensionarla y moverla. Una vez encontrada la posición deseada para el plano, pulsamos el botón ``\checkmark'' del panel de control para añadir el plano de forma definitiva.
\newpage
\imagenResize{0.7}{img/manualusuario/architect/floorpanel}{Pestaña ``Floors'' en el panel de control.}{floorpanel}

Para editar la planta basta con elegir el número de planta que quieres editar y en la pestaña ``Add'' de ``Floors'', subir un nuevo plano, que sobrescribirá al anterior.

Para borrar una planta, se debe elegir el número de planta y en la pestaña ``Delete'' de ``Floors'', confirmar el borrado.

\subsubsection{Añadir POIs}

Para añadir un poi debemos elegir un edificio con plano. Una vez elegido, en el apartado ``POI Toolbox'' de la pestaña ``POIs'' de nuestro panel de control, pinchamos en el icono de ``drag to add new POI'' y arrastramos a una posición del plano deseada.

\imagenResize{0.7}{img/manualusuario/architect/poispanel}{Pestaña ``POIs'' en el panel de control.}{poispanel}

Una vez asentado el poi, aparecerá un un panel, que emerge del icono del poi en el plano, donde se debe rellenar la información del poi (nombre del poi, descripción del poi y tipo de poi). Para añadirlo de forma definitiva se debe pulsar el botón ``Add'' del panel emergente.
\newpage
\imagenResize{0.7}{img/manualusuario/architect/poiiconplano}{Panel emergente del poi.}{poiiconplano}

Podemos encontrar otro tipo de poi, un \textit{conector poi}, que se coloca en el plano de la misma manera que el anterior, pero que se utiliza principalmente, para unir varios pois. Para seleccionar este poi en vez de pinchar y arrastrar desde el icono de ``drag to add new POI'', debemos pinchar y arrastrar desde el icono ``drag to add new connector'' .

\imagenResize{0.75}{img/manualusuario/architect/poiconector}{Panel emergente de conector poi.}{poiconector}

Para eliminar un poi basta con pinchar en el mismo y pulsar el botón con el texto ``x''.

Para crear un camino entre pois, debemos activar el ``toggle edge mode'' que se encuentra en el apartado ``POI Toolbox'' de la pestaña ``POIs'' de nuestro panel de control. Una vez activado debemos ir pinchando de poi en poi para crear un camino

\imagenResize{0.9}{img/manualusuario/architect/caminopoi}{Ejemplo de un camino creado}{caminopoi}.

Sabremos que el modo camino está activado si el botón de ``toggle edge mode'' está a on.  

\imagenResize{0.9}{img/manualusuario/architect/toggle}{``toggle edge mode'' on off.}{toggle}.

\subsubsection{Añadir ruta predefinida}

Para añadir una ruta predefinida a un edificio debemos activar ``add route mode'' que se encuentra en el apartado ``POI Toolbox'' de la pestaña ``POIs'' de nuestro panel de control.

\imagenResize{0.9}{img/manualusuario/architect/addmode}{``add route mode'' on off.}{addmode}


Una vez activado el modo añadir ruta, debemos pinchar en un poi que queramos añadir a la ruta, y aparecerá un desplegable donde nos encontramos un botón con el texto ``add POI to the route'', el cual debemos pulsar para añadir el poi a la ruta.

\imagenResize{0.75}{img/manualusuario/architect/poiaddroute}{Panel emergente del poi}{poiaddroute}

Cada vez que añadimos un poi a la ruta, se va coloreando la ruta. Una vez se hayan añadido todos los pois que se deseen, ponemos nombre a la ruta a través del campo ``Name'' de nuestro panel de control. Cuando se ha rellenado el campo con el nombre deseado, pulsamos el botón ``Create route'' para añadir la ruta de forma definitiva.  

\imagenResize{0.85}{img/manualusuario/architect/miruta}{Ejemplo de ruta coloreada.}{miruta}

Para editar la ruta, debemos poner el nombre de la ruta en el campo ``Name'' y añadir los pois que deseemos, lo que hará que se sobrescriba la ruta anterior.

Para visualizar las rutas debemos dirigirnos a la esquina superior izquierda de Architect y encontraremos un desplegable, que tiene 2 botones ``x'' para borrar la ruta seleccionada y ``Clear'' para no mostrar el trazado de la ruta. A la hora de seleccionar una ruta, aparece un número a la izquierda, que detalla la planta a la que pertenece la ruta. 

\imagenResize{0.85}{img/manualusuario/architect/selecruta}{Desplegable para seleccionar ruta, y botones de ``Clear'' y de borrado ``x''.}{selecruta}




\subsection{Viewer}

Para acceder a la herramienta Viewer, en la pagina de inicio de Geoindoor debemos pulsar el botón ``Viewer''. Viewer es la herramienta que utiliza un usuario para localizar puntos de interés de un edificio. Se visualizan edificios que han sido creados por los usuarios de Architect.
\newpage
\imagenResize{0.4}{img/manualusuario/viewer/viewerini}{Viewer inicio.}{viewerini}

\subsubsection{Búsqueda de edificio}

Para la búsqueda de un edificio se puede hacer de dos formas:

\begin{itemize}
	\item \textbf{Por nombre}: a través del desplegable de la parte superior.
	\item \textbf{Por posición}: a través del mapa e iconos.
\end{itemize} 

\imagenResize{0.8}{img/manualusuario/viewer/cluster}{Cluster de edificios.}{cluster}
\imagenResize{0.8}{img/manualusuario/viewer/gposicion}{Icono que indica la posición de un edificio.}{gposicion}

Para elegir un edificio basta con pulsar sobre los iconos de búsqueda por posición.

\subsubsection{Navegabilidad}

Una vez se ha elegido un edificio, se selecciona un poi. Al seleccionarlo aparece un panel emergente mediante el cual podemos elegir crear un camino automático (pulsando el botón de navegación) desde nuestra posición real (física), hasta el poi.

\imagenResize{0.8}{img/manualusuario/viewer/pemergente}{Panel emergente.}{pemergente}

\imagenResize{1}{img/manualusuario/viewer/btcamino}{Botón para crear el camino automático.}{btcamino}

\imagenResize{1}{img/manualusuario/viewer/caminocreado}{Camino automático.}{caminocreado}

\newpage
Si miramos a la derecha de interfaz, veremos unas flechas que sirven para elegir que planta mostrar del edificio.

\imagenResize{1}{img/manualusuario/viewer/planta}{Selector de planta del edificio.}{planta}

En la esquina inferior izquierda, encontramos unos botones que nos sirven para mostrarnos nuestra posición física, ajustar nuestra posición física, y mostrar la posición del poi seleccionado.

\imagenResize{1}{img/manualusuario/viewer/btnav}{Botones de navegación.}{btnav}

\imagenResize{1}{img/manualusuario/viewer/mpoi}{Botón para mostrar la posición del poi.}{mpoi}
\imagenResize{1}{img/manualusuario/viewer/mposicion}{Botón de mostrar posición del usuario.}{mposicion}
\imagenResize{1}{img/manualusuario/viewer/ajpos}{Botón de ajuste de posición.}{ajpos}

\subsubsection{Rutas predefinidas}

Para mostrar las rutas predefinidas de un edificio nos dirigimos a la parte superior izquierda, y encontramos un desplegable con las rutas predefinidas del edificio. A la izquierda del nombre de las rutas se indica la planta a la que pertenece (0, 1, 2 \ldots). Cada vez que se selecciona una ruta, se dibuja su traza en el plano. Para dejar de mostrar la ruta basta con pulsar el botón ``x''.

\imagenResize{1}{img/manualusuario/viewer/despleg}{Desplegable para la selección de rutas.}{despleg}

\imagenResize{1}{img/manualusuario/viewer/rutadib}{Dibujado de la ruta seleccionada.}{rutadib}



\bibliographystyle{plain}
\bibliography{bibliografiaAnexos}

\end{document}
